\begin{abstract}
The rapid growth of task-specific deep neural networks has popularized weight-space model merging as a zero-overhead alternative to classic Mixture-of-Experts (MoE) or ensembling. Recently, several works have asserted a theoretical ``Layer-Averaging Collapse'' (or rank-1 collapse), claiming that dynamic layer-wise routing coefficients are redundant and inevitably collapse to a single global trajectory. In this paper, we deconstruct this claim from a critical methodological perspective. We show that the rank-1 collapse theorem is a pure artifact of over-simplified, linear representation-space sandboxes and low-conflict multi-task environments. By conducting a rigorous, multi-seed physical empirical evaluation on Split-MNIST digits subsets using deep neural network backbones (DeepMLP-12 and TinyCNN-4) across task suites of varying semantic conflict, we expose the limits of representational collinearity. Singular Value Decomposition (SVD) of the learned physical layer-wise coefficient matrices reveals that the Collinearity Ratio drops to $0.4987 \pm 0.08$ (DeepMLP-12) and $0.5673 \pm 0.03$ (TinyCNN-4) under Cross-Domain task conflict, proving that depth-specialized routing policies emerge in real physical architectures as a semantic necessity. Furthermore, our results clarify the parameter-variance trade-offs of Dynamic Routing under few-shot budgets, suggesting that the community must move away from toy linear sandboxes and adopt more rigorous physical evaluation protocols to properly understand and leverage weight-space dynamic model merging.
\end{abstract}
